\section*{附录 A\quad 数据字典、尺度核验与样本对齐明细}
\addcontentsline{toc}{section}{附录 A\quad 数据字典、尺度核验与样本对齐明细}

本附录列出附件一 231 个特征的名称、尺度判定（校准对数、校准线性、未校准信号、综合指标）与单位口径，250 款啤酒的样本对齐（行号与编号对照）与附件一和附件八的冲突单元明细，以及附件二、附件六的复算误差逐项结果。

\section*{附录 B\quad 50 维完整评分、等级概率与逐维评价}
\addcontentsline{toc}{section}{附录 B\quad 50 维完整评分、等级概率与逐维评价}

本附录转置给出 50 行乘 12 款跨日期品评啤酒的完整评分表，单元写为预测等级 [90\% 区间] / 实测等级；随后给出 250 款乘 50 维的逐维预测等级、等级概率与 90\% 等级区间作为机器可读补充文件，以及逐维连续 $R^2$、等级 MAE、QWK 与半组信度。

\section*{附录 C\quad 前三项关键化合物、稳定性与化学证据完整表}
\addcontentsline{toc}{section}{附录 C\quad 前三项关键化合物、稳定性与化学证据完整表}

本附录按六个维度组给出 49 维乘 3 名的 147 项完整表，含选择得分、95\% 置信区间、风格内系数、来源标签、气味活度值占比与模块编号，并列出来源与缺陷标签的逐项证据。

\section*{附录 D\quad 各风格生成谱、感官画像与迁移明细}
\addcontentsline{toc}{section}{附录 D\quad 各风格生成谱、感官画像与迁移明细}

本附录给出 22 个风格各 1,000 个生成化合物谱的摘要统计、50 维感官画像区间，以及 Pils/Lager 到 Tripel 的 11 款迁移明细与全部有序风格对的迁移汇总。

\section*{附录 E\quad 参数网格、补充对照与敏感性分析完整结果}
\addcontentsline{toc}{section}{附录 E\quad 参数网格、补充对照与敏感性分析完整结果}

本附录列出问题一、二、三的实验设置、完整搜索网格、各外层折所选参数的频数、全部对照模型与敏感性臂的逐项结果。

\begin{table}[H]
\centering\small
\caption{问题一实验设置与对照}
\label{tab:q1-setup}
\begin{tabular}{p{3.0cm}p{5.4cm}p{4.2cm}}
\toprule
环节 & 实验设置 & 控制或检验 \\
\midrule
目标与输入 & 250 款啤酒、231 个化学特征、50 个感官维度；以逐品评员原始评分的每款每维均分为连续目标，均分取整为 0--6 级真值 & 预测目标、特征与真值同口径 \\
\addlinespace
化学缺失 & 空值在训练折内按列均值填充 & 填充不泄漏测试信息 \\
\addlinespace
数据划分 & 根据日期和人员确定 25 个日期组，按日期组隔离、按风格分层的 5 折 $\times$ 5 次重复，种子 0 到 4，内层同为 4 折 & 同日品评信息不泄漏，调参与报告分开 \\
\addlinespace
连续模型 & 每维一个平方损失 HistGBR，六组内共享超参数，内层按组内平均 RMSE/SD 选参 & 控制模型复杂度 \\
\addlinespace
搜索范围 & 深度 $\{2,3\}$、叶节点 $\{10,20\}$、学习率 0.05、迭代 $\{100,200,400\}$、$l_2$ 正则 $\{0,1\}$、六个单调约束开关由内层选择 & 选参可复现，正文另述迭代下边界检查 \\
\addlinespace
等级与区间 & 内层折外残差构造 0--6 等级概率；分布中位数为主输出，分位数切点为第二输出；90\% 等级区间 & 连续预测转为可交付等级及其不确定性 \\
\addlinespace
对照与判据 & 风格收缩均值为基线，ElasticNet 检查线性替代，随机森林，只用附件四的两阶段方法对照；主指标等级 MAE，辅助 QWK、一致率、$\pm1$ 内比例；按日期组 bootstrap 2,000 次给配对差 95\% 置信区间 & 每个实验为何存在，什么结果支持主张 \\
\bottomrule
\end{tabular}
\end{table}

\begin{table}[H]
\centering\small
\caption{问题二实验设置与对照}
\label{tab:q2-setup}
\begin{tabular}{p{3.0cm}p{5.4cm}p{4.2cm}}
\toprule
环节 & 实验设置 & 控制或检验内容 \\
\midrule
输入与目标 & 250 款啤酒；附件一 231 列剔除 5 个综合指标与内标，剩 225 个候选；49 个维度，barnyard 记作同 stale\_hops；感官目标为原始评分均分的秩 & 候选都是化合物或理化指标，目标与问题一同口径 \\
\addlinespace
缺失与秩 & 8 个单化合物列共 60 个空值按列均值填充，每次 bootstrap 样本内重算；信号下限按并列最低秩处理 & 填充值不解释为浓度 \\
\addlinespace
图与模块 & 边权取 $\lvert\rho\rvert$，$1-\lvert\rho\rvert$ 上平均连接层次聚类，$\tau=0.6$ 切割 & 模块内平均相关至少约为 $\tau$ \\
\addlinespace
选择规则 & 第 1 位取最大 $\lvert r_{jd}\rvert$，第 2、3 位取控制已选集合后的秩偏相关最大者，同一模块至多入选 1 个 & 用条件关联去除间接相关者 \\
\addlinespace
显著性 & 置换 1,000 次，第 1 位用全体候选最大相关零分布的 95\% 分位，第 2、3 位用 Freedman--Lane 置换，并同时要求日期组 bootstrap 95\% 置信区间不含 0 & 维度内多重比较与日期组内相关 \\
\addlinespace
稳定性 & 25 个日期组有放回重抽 $B=1000$，每次重建相关矩阵与模块，报告化合物层与模块层入选频率 $\pi^c$、$\pi^m$，分档界限 0.6 与 $1/3$ & 结论对样本扰动的稳定程度 \\
\addlinespace
规则依据实验 & 真实化学矩阵植入 3 个不同模块的化合物，权重 (1,0.75,0.5) 与 (1,1,1)，第 1 个植入化合物的 $r$ 取 0.2、0.3、0.45，每种设置 1,000 次，噪声含日期组偏移 & 只用化学矩阵、不用感官数据，回应根据结果改规则的质疑 \\
\addlinespace
对照方法 & 无图约束边际选择、图约束边际选择、ElasticNet、问题一梯度提升树、风格内分析 & 每种对照对应方法的一处改动 \\
\bottomrule
\end{tabular}
\end{table}

\begin{table}[H]
\centering\small
\caption{问题三实验设置与对照}
\label{tab:q3-setup}
\begin{tabular}{p{3.0cm}p{5.4cm}p{4.2cm}}
\toprule
环节 & 正文报告的设置 & 控制或检验内容 \\
\midrule
数据与生成列 & 250 款、22 种细分风格、225 个生成列；综合指标由成分派生，内标取训练中位数 & 生成对象是单体化合物 \\
\addlinespace
表示 & 训练折经验分布映入正态得分；共同下限为并列原子，取截断正态条件期望 & 不同尺度与未检出状态统一处理 \\
\addlinespace
模型与选参 & 风格均值收缩 $\kappa\in\{0,1,2,5,10,20\}$，因子数 $k\in\{2,5,10,15,20,25,30\}$，纯因子与因子加模块块两种结构，对角与 Ledoit--Wolf 作对照；外层训练集内 5 折按留出潜空间评分选参 & 选参与报告分离，风格内共变的结构 \\
\addlinespace
数据划分 & 按风格分层、以单款啤酒为单位的 5 折 $\times$ 5 次重复；画像用问题一的日期组外层划分 & 化学测量与品评日期无关，画像防止同日评价信息泄漏 \\
\addlinespace
评价器 & 随机森林 500 棵树只用真实啤酒训练；C2ST 置换 200 次并做真实对真实检验；按啤酒 bootstrap 2,000 次 & 判别器与生成器不共享高斯假设，检验能力由阳性对照给出 \\
\addlinespace
对照生成器 & 纯因子、无风格、逐列独立、复制、风格内插值五个对照 & 检验风格条件、化合物共变、新颖性 \\
\addlinespace
画像 & 每个风格生成 1,000 个化合物谱，经问题一全数据模型与残差池得到等级分布 & 画像只由化学谱决定 \\
\addlinespace
迁移约束 & 化学层 9 列，阈值取附件八均值口径；感官层 6 个缺陷维度预测等级不升高；沿路径 $\alpha$ 取 0 到 1、步长 0.05，原尺度复查 & 把任务的不恶化写成可检验的条件 \\
\addlinespace
迁移对照 & 平移、截断、原尺度平移三个消融，原谱不变、截断加感官搜索、检索真实目标风格三个同约束对照 & 约束、投影与潜空间各自的作用 \\
\bottomrule
\end{tabular}
\end{table}

\section*{附录 F\quad 软件环境、运行命令与全部源程序}
\addcontentsline{toc}{section}{附录 F\quad 软件环境、运行命令与全部源程序}

本附录给出求解与绘图所用的软件版本、随机种子、运行命令与全部源程序清单。
